\section{A Structural Model of Linguistic and Semantic Components}
\label{sec:model}

Motivated by the empirical observations in Section~2, we introduce a structural
model that captures the essential geometric properties of language model
representations relevant to anisotropy and spike suppression.
The purpose of this model is not to reproduce all details of neural
representations, but to isolate the minimal structure needed to explain the
empirical behavior of averaging observed later in this section.
A formal explanation of these behaviors is provided in Section~4.

\subsection{Token Representation Model}
\label{subsec:token_model}

Let $x_i \in \mathbb{R}^H$ denote the representation of token $i$ at a fixed layer,
and let
\[
X = 
\begin{bmatrix}
x_1^\top \\
\vdots \\
x_N^\top
\end{bmatrix}
\in \mathbb{R}^{N \times H}
\]
collect the representations of all tokens in a batch.
We model each token representation as the sum of two components,
\begin{equation}
\label{eq:model_decomposition}
x_i = \ell + s_i,
\end{equation}
where:
\begin{itemize}
  \item $\ell \in \mathbb{R}^H$ is a \emph{shared linguistic component}, common to
  all tokens in the batch; and
  \item $s_i \in \mathbb{R}^H$ is a \emph{semantic component} that varies across
  tokens.
\end{itemize}

This decomposition reflects the empirical findings in
Sections~\ref{subsec:spike_long_tail}--\ref{subsec:tail_incoherent}:
linguistic structure (e.g., grammar, syntax, and shared functional patterns) is
globally coherent across tokens, while semantic content is diverse and
context-dependent.

\subsection{Structural Assumptions}
\label{subsec:structural_assumptions}

The empirical analyses in Sections~2.2--2.4 motivate the following structural
assumptions.

\paragraph{Assumption 1 (Low-rank linguistic structure).}
The shared component $\ell$ lies in a low-dimensional subspace
$\mathcal{L} \subset \mathbb{R}^H$ of dimension $d_\ell \ll H$.
Equivalently, $\ell$ is well approximated by the span of a small number of leading
right singular vectors of $X$, consistent with the spike structure observed in
Figure~\ref{fig:spectrum_spike_tail_combined} and the coherence analysis in
Figure~\ref{fig:colmean_alignment}.

\paragraph{Assumption 2 (Semantic compressibility).}
Each semantic component $s_i$ lies predominantly in the orthogonal complement
$\mathcal{L}^\perp$ and is \emph{compressible} in a fixed semantic eigenbasis:
most of its energy is concentrated on a subset of coordinates whose size is small
relative to $\dim(\mathcal{L}^\perp)$.
This assumption is empirically supported by the effective support statistics
shown in Figure~\ref{fig:tail_effective_support}.

\paragraph{Assumption 3 (Semantic incoherence across sequences).}
For tokens drawn from different sequences, or from sufficiently distant
positions within a long sequence, semantic components are weakly correlated in a
normalized sense.
Let $s_i^{\perp} = P_{\mathcal{L}^\perp}s_i$ denote the projection of $s_i$ onto the
semantic subspace.
Then for $i \neq j$ from different sequences,
\begin{equation}
\label{eq:assump_cos_incoh}
\mathbb{E}\!\left[
\frac{\langle s_i^{\perp},\, s_j^{\perp}\rangle}
{\|s_i^{\perp}\|_2\,\|s_j^{\perp}\|_2}
\right]
\approx 0,
\end{equation}
with bounded second moments $\mathbb{E}\|s_i\|_2^2 \le \sigma^2$.
This assumption is directly validated by the coefficient-space cosine similarity
distributions in Figure~\ref{fig:tail_coeff_cosine}.

\subsection{Empirical Effect of Averaging on Spike and Tail}
\label{subsec:avg_empirical}

We now empirically examine how mean removal (averaging) affects the spectral
structure characterized in Section~2.
The goal of this subsection is purely observational: to quantify how much of the
spectral spike is suppressed by averaging, and how much of the long tail is
preserved, without appealing to any theoretical explanation.
A formal account of these phenomena is deferred to Section~4.

\paragraph{Experimental setup.}
Given the token matrix $X \in \mathbb{R}^{N\times H}$, we consider two operations:
(i) \emph{mean removal},
\[
X_{\mathrm{avg}} = X - \mathbf{1}\mu^\top,
\quad
\mu = \tfrac{1}{N}\sum_{i=1}^N x_i,
\]
and (ii) \emph{top-$k$ SVD removal},
\[
X_k = X - X P_k,
\]
where $P_k$ projects onto the top-$k$ right singular vectors of $X$.
We compare these operations using Frobenius-norm energy ratios, which measure how
much variance remains after each transformation.

\paragraph{Mean alignment with the spike subspace.}
Figure~\ref{fig:avg_mean_in_spike} reports the fraction of mean energy captured by
the spike subspace, $\|P_k\mu\|_2^2 / \|\mu\|_2^2$, as a function of $k$.
The batch mean lies almost entirely within the spike subspace for modest $k$; for
example, projecting onto the top $1\%$ of singular directions captures over
$99.9\%$ of the mean energy.
This indicates that mean removal primarily targets directions aligned with the
dominant spectral spike.

\paragraph{Spike energy reduction and tail preservation.}
Figures~\ref{fig:avg_spike_reduction} and~\ref{fig:avg_tail_preservation} quantify
the effect of averaging on the spike and tail components, respectively.
Mean removal suppresses a nontrivial portion of spike energy, while preserving
nearly all tail energy across a wide range of spike-subspace definitions.
In contrast, top-$k$ SVD removal progressively eliminates both spike and tail
energy as $k$ increases.

\paragraph{Conservative nature of averaging.}
Figure~\ref{fig:avg_residual_energy} compares the total residual energy after mean
removal and after top-$k$ SVD removal.
Mean removal eliminates a fixed shared component, resulting in a bounded energy
reduction, whereas SVD removal continues to remove energy as the projection rank
grows.
This highlights the conservative nature of averaging: it suppresses the dominant
shared structure without aggressively removing variance from the long tail.

\begin{figure}[t]
  \centering
  \begin{subfigure}[t]{0.48\linewidth}
    \centering
    \includegraphics[width=\linewidth]{figures/avg_mean_in_spike_vs_k.png}
    \caption{Mean energy captured by the spike subspace.}
    \label{fig:avg_mean_in_spike}
  \end{subfigure}
  \hfill
  \begin{subfigure}[t]{0.48\linewidth}
    \centering
    \includegraphics[width=\linewidth]{figures/avg_spike_reduction_vs_k.png}
    \caption{Spike energy reduction under mean removal.}
    \label{fig:avg_spike_reduction}
  \end{subfigure}

  \vspace{0.5em}

  \begin{subfigure}[t]{0.48\linewidth}
    \centering
    \includegraphics[width=\linewidth]{figures/avg_tail_preservation_vs_k.png}
    \caption{Tail energy preservation under mean removal.}
    \label{fig:avg_tail_preservation}
  \end{subfigure}
  \hfill
  \begin{subfigure}[t]{0.48\linewidth}
    \centering
    \includegraphics[width=\linewidth]{figures/avg_residual_energy_comparison.png}
    \caption{Residual energy: mean removal vs.\ top-$k$ SVD removal.}
    \label{fig:avg_residual_energy}
  \end{subfigure}

  \caption{Empirical validation of averaging as conservative spike suppression.
  Mean removal targets directions aligned with the spike subspace
  (top-left), suppresses spike energy (top-right), preserves long-tail energy
  (bottom-left), and removes less total energy than top-$k$ SVD projection
  (bottom-right).}
  \label{fig:avg_vs_svd_validation}
\end{figure}

\paragraph{Summary of empirical findings.}
Across layers and batches, we consistently observe that mean removal:
(i) targets directions aligned with the spectral spike;
(ii) suppresses a substantial fraction of spike energy; and
(iii) preserves the vast majority of long-tail energy.
These empirical regularities set the stage for the next section, where we
provide a formal explanation for why averaging exhibits these properties under
the structural assumptions validated above.